\documentclass[12pt]{article}
\usepackage[utf8]{inputenc}
\usepackage[T1]{fontenc}
\usepackage{hyperref}
\title{Tokharian B Language Resources for K-12 Curriculum Development}
\author{Compiled by Claude Sonnet 5 (claude-sonnet-5) for Andr\'as Kornai}
\date{}
\begin{document}
\maketitle

\section{Introduction}
Tokharian (Tocharian) B is an extinct Indo-European language of the Tarim Basin, attested in manuscripts roughly of the 5th to 10th centuries AD (sources differ on the dating). It has no native speakers and no known community of heritage learners, and no K-12 programme was found. All resources below are scholarly or university-level; any school curriculum would have to be built from them and adapted heavily.

\section{Primary Educational Resources}
\begin{itemize}
\item Tocharian Online (University of Texas at Austin, Linguistics Research Center): lessons, texts from Tocharian A and B, and a glossary. \url{https://lrc.la.utexas.edu/eieol/tokol/200}
\item Michael Weiss, \emph{Kuś\'i\~n\~ne Kantwo: Elementary Lessons in Tocharian B}, described as the first pedagogical grammar of Tocharian B. Announcement: \url{https://linguistics.cornell.edu/news/new-book-helps-students-learn-ancient-tocharian-language}
\end{itemize}

\section{Linguistic and Reference Materials}
\begin{itemize}
\item Douglas Q. Adams, \emph{A Dictionary of Tocharian B} (revised and enlarged edition, Rodopi, 2013); review: \url{https://journals.phil.muni.cz/linguistica-brunensia/article/view/30841}
\item Etymological Dictionary of Tocharian B (University of Gothenburg database): \url{https://databaser.ub.gu.se/etymological-dictionary-of-tocharian-b/188896}
\item CLTK documentation page for Tocharian B (includes a Swadesh list): \url{https://cltk.readthedocs.io/en/latest/tocharian_b.html}
\item Starling database Tocharian wordlist (PDF): \url{https://starlingdb.org/new100/txr.pdf}
\end{itemize}

\section{Religious and Community Resources}
The surviving texts are mainly Buddhist, but no modern community or religious organisation using the language was found in the search results.

\section{Digital and Audio Resources}
Tocharian Online (above) and the CLTK page are the digital resources found. No audio recordings were found. A third-party PDF copy of Adams's dictionary appeared in results, but its provenance is doubtful and it is not cited here.

\section{Curriculum Development Context}
There is no living speech community, so a K-12 unit could at most be an enrichment module (writing systems, historical linguistics, Silk Road history) using simplified glossary items from the resources above. Teachers would need to prepare original materials, and the content should be vetted by a specialist.

\end{document}